\section{Introduction}
Salivary gland mucosa-associated lymphoid tissue (MALT) lymphoma and Sjögren's syndrome (SS)-associated inflammation can share lymphoid-rich histological appearances. SS is also a clinical setting for lymphomagenesis \citep{ref1}. Lymphoepithelial lesions and architectural changes occur in both processes, making the composition and distribution of infiltrates relevant to pathological assessment \citep{ref2}. The diagnostic distinction rests on integrated morphology and, when indicated, immunophenotypic and molecular findings within current lymphoma classification frameworks \citep{ref3,ref4}. An image-based classifier could provide an additional signal for review, but its diagnostic scope must be established in the relevant case spectrum.

Computational studies have addressed SS diagnosis using cell--tissue graphs and automated evaluation of minor salivary gland focus scores \citep{ref10,ref23}. Related lymphoma studies have distinguished orbital MALT lymphoma from IgG4-related ophthalmic disease or classified MALT and diffuse large B-cell lymphoma in selected histological material \citep{ref11,ref24}. These studies differ in tissue, reference categories and image-sampling strategy. They motivate direct evaluation of salivary gland MALT-versus-SS classification without making their reported accuracy interchangeable with this task. Here, MALT and SS denote the recorded slide-level diagnostic groups; an SS medical history in a patient with MALT is a separate clinical attribute.

Weakly supervised multiple-instance learning (MIL) can learn slide-level classifications without exhaustive lesion annotations \citep{ref5,ref6,ref7}. Pathology foundation models provide reusable patch representations, allowing disease-specific aggregation to be trained from cached embeddings \citep{ref8,ref9}. Recent benchmarking and the nnMIL framework emphasize that feature representations and training recipes must be considered alongside aggregation architecture \citep{ref25,ref26}. Simple feature-based controls and comparisons of complete implementations therefore provide useful context for evaluating a complex model on a small disease-specific cohort.

We studied retrospective MALT--SS classification using fixed UNI2-h embeddings and a custom Hybrid-Evidence MIL aggregator. The model combines attention, distributional and prototype-margin summaries to produce a slide-level score. We evaluated its discrimination against alternative MIL implementations, assessed the selected checkpoint on an external collection and examined model-selected regions. Because internal validation slides also selected checkpoints and the external material has mixed archive stain tags, we distinguish development estimates from independent clinical validation. Additional feature-based controls and implementation analyses accompany the primary results.
